% Programming with InitechBase, Chapter 2: Program Structure and Flow
% (AN-16; annex Sec 8.2 P2).  Panels replayed (examples/p2-*.dot); sources
% cmd/flow.c (IF/ELSE/ENDIF, DO WHILE/ENDDO/LOOP/EXIT, DO CASE; guards must
% be Logical, #37).
\chapter{Program Structure and Flow}
\Gates{test-pubs-examples, test-interp-flow}
A program runs its commands in order, from the first line to the last, unless a
structure tells it otherwise.  InitechBase has three: \cmd{IF} chooses whether
to carry out a group of commands, \cmd{DO CASE} chooses one group among
several, and \cmd{DO WHILE} repeats a group.  This chapter describes each, with
the programs of the system that use them.

\begin{ChapterCovers}
\item Sequence
\item \cmd{IF}, \cmd{ELSE} and \cmd{ENDIF}
\item \cmd{DO CASE}
\item \cmd{DO WHILE}, \cmd{LOOP} and \cmd{EXIT}
\item Nesting and indentation
\item Conditions
\end{ChapterCovers}

\ATEHead{Preparing for This Chapter}
Start InitechBase with a fresh copy of the sample diskette, and open the ledger
with \type{USE INVOICE}.

\ATEHead{Sequence}
Without a structure, the commands of a program run one after the other, once
each.  The first five commands of \file{OPENINV.PRG} (Chapter~1) are a sequence:
two headings, a memory variable and a move to the first record.

\ATEHead{Decisions: IF}
\cmd{IF} is followed by a condition.  When the condition is true, the commands
up to \cmd{ELSE}, or up to \cmd{ENDIF} if there is no \cmd{ELSE}, are carried
out.  When it is false, the commands after \cmd{ELSE} are carried out instead.
Either way, the program goes on after \cmd{ENDIF}.

\SyntaxLine{\kw{IF} \xp{condition}}
\SyntaxLine{\ \ \xp{commands when true}}
\SyntaxLine{\kw{ELSE}}
\SyntaxLine{\ \ \xp{commands when false}}
\SyntaxLine{\kw{ENDIF}}

\file{INVTOTAL.PRG} adds every amount to the total of the ledger, and the paid
amounts to a second total:
\ProgramListing{INVTOTAL.PRG}

\ATEHead{Conditions}
A condition must be a logical value: a logical field such as \texttt{PAID}, a
comparison, a logical function such as \fn{EOF()} or \fn{FOUND()}, or these
joined with \texttt{.AND.}, \texttt{.OR.} and \texttt{.NOT.}.  InitechBase does
not treat a number or a character string as true or false.  A program whose
condition is of another type stops with
\xmsg{37}{Not a Logical expression.}

\ATEHead{Choices: DO CASE}
\cmd{DO CASE} holds several \cmd{CASE} lines, each with a condition.  The
commands of the first \cmd{CASE} whose condition is true are carried out, and
the program goes on after \cmd{ENDCASE}.  If no condition is true, the commands
after \cmd{OTHERWISE} are carried out.

\SyntaxLine{\kw{DO CASE}}
\SyntaxLine{\kw{CASE} \xp{condition 1}}
\SyntaxLine{\ \ \xp{commands}}
\SyntaxLine{\kw{CASE} \xp{condition 2}}
\SyntaxLine{\ \ \xp{commands}}
\SyntaxLine{\kw{OTHERWISE}}
\SyntaxLine{\ \ \xp{commands}}
\SyntaxLine{\kw{ENDCASE}}

The main menu of the system (Chapter~11) carries out the task the clerk
chooses with a \cmd{DO CASE}, one \cmd{CASE} for each choice.

\ATEHead{Counting by Class}
\file{DUECLASS.PRG} counts the invoices of the ledger in three classes of
amount.  The \cmd{DO CASE} inside the loop chooses the count for each record:
\ProgramListing{DUECLASS.PRG}
\DotExample{examples/p2-dueclass}
The order of the \cmd{CASE} lines matters.  An amount under 100 is also under
1000, but only the first \cmd{CASE} that is true is carried out, so it is
counted once.

\ATEHead{Loops: DO WHILE}
\cmd{DO WHILE} tests its condition; if it is true, the commands up to
\cmd{ENDDO} are carried out and the condition is tested again.  The loop ends
when the condition is false.

\SyntaxLine{\kw{DO WHILE} \xp{condition}}
\SyntaxLine{\ \ \xp{commands}}
\SyntaxLine{\kw{ENDDO}}

The loop of nearly every program of the system processes the records of the
ledger:

\SyntaxLine{\kw{GO TOP}}
\SyntaxLine{\kw{DO WHILE .NOT. EOF()}}
\SyntaxLine{\ \ \xp{commands for one record}}
\SyntaxLine{\ \ \kw{SKIP}}
\SyntaxLine{\kw{ENDDO}}

\cmd{LOOP} goes back to the test at once, leaving out the rest of the loop;
\cmd{EXIT} leaves the loop at once.  The menu of Chapter~11 is a loop that
never ends by its condition, \cmd{DO WHILE .T.}: \cmd{EXIT} ends it when the
clerk chooses to leave.

The memory variables of the program remain after it ends, so you can look at
them at the dot prompt:
\DotExample{examples/p2-openinv}

\ATEHead{Leaving a Loop Early}
\file{FIRSTBIG.PRG} looks for the first open invoice of 1000 or more.  It
passes over paid invoices with \cmd{LOOP} -- note the \cmd{SKIP} before it, or
the loop would test the same record for ever -- and leaves the loop with
\cmd{EXIT} as soon as it finds one:
\ProgramListing{FIRSTBIG.PRG}
\DotExample{examples/p2-firstbig}
The logical memory variable \texttt{MFOUND} remembers why the loop ended: at
the end of the file, or at an invoice found.

\ATEHead{Nesting and Indentation}
A structure may hold other structures: an \cmd{IF} inside a \cmd{DO WHILE}, as
in \file{OPENINV.PRG}, or a \cmd{DO CASE} inside an \cmd{IF}.  Each must be
closed inside the one that holds it.  InitechBase ignores blanks at the
beginning of a line, so indent the commands inside each structure by two
blanks: the programs of the system are written that way, and the structure is
then plain to see.

\begin{Tip}
Write the closing line of a structure -- \cmd{ENDIF}, \cmd{ENDCASE},
\cmd{ENDDO} -- as soon as you write its first line, then fill in the commands
between.
\end{Tip}

\ATEHead{Where Do You Go From Here?}
The programs of this chapter keep their totals in memory variables.  Chapter~3
tells you how memory variables are created and how long they last.

\ATEHead{Summary}
\begin{enumerate}
\item Commands run in sequence unless a structure decides or repeats.
\item \cmd{IF} \ldots{} \cmd{ELSE} \ldots{} \cmd{ENDIF} chooses between two groups.
\item \cmd{DO CASE} \ldots{} \cmd{ENDCASE} chooses the first group whose condition
  is true.
\item \cmd{DO WHILE} \ldots{} \cmd{ENDDO} repeats; \cmd{LOOP} and \cmd{EXIT} jump.
\item Every condition is a logical value.
\end{enumerate}
